\chapter{Full stack web design}\label{fullstack}

This section will introduce the concept of the tech stack and discuss concepts of full-stack web development and their relevance to this project's goals. Starting with a primer on the basics of web applications and information transfer over the internet, following with sections on conceptual differences between single- and multi-page applications, concepts and examples in the software architecture of web apps, continuing with an overview of the deployment process and the role of containerization, and concluding with an assessment of the process behind constructing a tech stack for this project.

\section{The Tech Stack in Web Development}

In the field of computer programming, a software stack is the set of software components required to run and support an application \cite{adedeji_full-stack_2023}. The stacks one uses for different applications can vary, but popular combinations of components have proliferated over time due to the compatibility of certain components with one another. In web development, a common classification of stack components is labeling tools as part of the "frontend", which encompasses the presentation and organization of a site's content, determining how a site is made to appear to a user; or the "backend", which encompasses the logic of the site, processing user requests by routing them to the right component, storing data necessary to the function of the software, and managing application operations. Some also use "middleware" as a separate term, covering steps like authentication, the routing of information between components, and liaising with third-party services. Going forward, this paper will use the term "middleware" in this way.

An example of a popular web development stack is MEAN, an acronym standing for the four components it contains:

\begin{itemize}
	\item \textbf{MongoDB}, A database management system to store information needed by the application. It is non-relational, a concept explored in \autoref{databases}
	\item \textbf{Express.js}, a framework for designing web applications with JavaScript
	\item \textbf{Angular}, a framework for designing user-interfaces
	\item \textbf{Node.js}, a runtime environment allowing JavaScript to be executed on a server, outside of a browser.
\end{itemize}

Each component in the stack fulfills a different purpose, and together are able to run a web application. Node.js and MongoDB make up the backend, performing the processing of information and storing data, respectively. Angular handles the frontend, providing the tools to build a flexible user-interface. Finally, Express.js is the middleware -- it routes code and information between the other modules. Since MEAN does not require a specific operating system to run within \cite{noauthor_what_nodate}, it is known as an "OS-agnostic" stack. An example of the opposite -- an "OS-level" stack -- is LAMP, which uses the Linux operating system. The rest of the LAMP acronym stands for Apache, the framework to run the backend server; MySQL, the database manager; and Perl, as the primary coding language used.

Even when part of the same unifying software, different stack components must still have an established method of communicating with one another. As discussed previously in \autoref{bg}, components communicate with one another through APIs. HTTP (usually through a RESTful API) forms the boundary between the frontend and backend, while middleware and backend components have their own methods of communicating with one another. Examples of how multiple components can be organized within a system will be discussed later in the chapter.

There is no set framework for what kinds of component constitute a specific "stack"; rather, named stacks like LAMP and MEAN highlight the most important and relevant components. Another factor is that certain tools may overlap with the functionality of others (for example, one backend framework may include functionality to communicate directly with a database while another may need an intermediary), meaning a rigid categorization of components would not be helpful. Types of component that exist in one stack might be redundant or unnecessary in another.

\section{Single page vs. Multipage}

One key design decision crucial in the early stages of creating web applications is the choice between creating a single-page or multi-page application. The distinction is fairly intuitive: single-page applications load only one HTML page initially and use client-side code like JavaScript to perform dynamic functionality such as fetching data and changing the user-interface. Multi-page applications, on the other hand, consist of more than one HTML page \cite{adedeji_full-stack_2023}. While the actual size of the sitemap in these applications can vary widely -- anywhere from a small handful of pages to dozens -- it is seen as a significant distinction to go from one page to multiple, no matter the exact level of the difference.

The key difference marking the distinction between the two types is that in single-page applications, when a user first lands on it, the entire frontend of the application is downloaded by their browser. Even if the application has multiple pages (which, while profoundly unintuitive, does not preclude it from being called "single-page"), the content of each page is sent from the server to the client in one response, and any page-to-page routing is handled by the application's client-side code that the user's browser executes. If, through the user's use of the application, more information is needed from the server, an API call will be made to request it. But the overall distinction between single-page and multi-page is how the application's content is delivered to the user.

Single-page applications are a relatively new concept in the history of the web. In the early days of the internet, webpages would each be sent from server to client in discrete requests, much like a multi-page application today. Over time, as personal computers and web browsers became more advanced, the capability of users' devices to provide interaction and functionality on the client-side increased through tools such as JavaScript. Over time, client-side functionality advanced to the point where for certain web tools it made more sense to send a website's entire functionality to the user all at once and have it act more like a desktop application than a traditional website. Most modern websites use some level of client-side functionality, whether they're single- or multi-page, but a trade-off between the two still exists: for some sites it's beneficial to send the user a full application all at once, and for others it would be more efficient to separate different sections of the site's functionality and only need to deliver the content the user actually needs.

The decision whether to structure an application as single-page or multi-page depends on how a designer wants information to be organized. To give an example of each, we can look at Amazon and Google Maps. Amazon is a multi-page application where information about each product is contained on its own page. Because its catalog of items is so vast, it makes sense to cordon off each individual item into its own area with relevant information like price, availability, and user reviews. Other areas of the site like the search results page and a user's shopping cart are also their own pages, only being fetched and delivered to the user if they specifically ask (although a smaller sidebar displaying items in a user's cart is shown while browsing, the cart page contains more details not included on it).

Google Maps, on the other hand, mostly operates on one homepage, displaying a basemap of the globe, focused in on a user's location. The site's main features are displaying information about a selected location and mapping a route of transport between two places -- both of which are accomplished without leaving the main page. Users select a place either by clicking or with the search bar. Information about the place is then displayed, along with options to give directions, save it to a list, etc. Despite its single-page nature, however, the entirety of information Google contains about world locations is -- as one would expect -- not saved into a user's browser as soon as they land on Google Maps. Rather, an API call is made when a user selects a location, requesting further information from Google's servers such as images and user reviews if the place is an established business. Despite needing additional information from the backend, the site incorporates the retrieved information into the existing page rather than navigating to a separate page for each location as Amazon does with its products.

\section{Architecture}

The tech stacks in modern software and development often comprise multiple components that interact to perform its functionality. As we saw earlier, different stacks can contain or emphasize different types of components while being agnostic on others. The structure of these components and how they connect with one another can be used as a basic definition of software architecture. Depending on the needs of the software for its stakeholders, it may need a different combination of components, or commonly-used components assembled and connected in a different way \cite{ghidersa_software_2022}. In web development, while the frontend and backend are commonly-used categorizations of services and components, there are many ways to structure an application, some blurring the lines between the categories.

This section will cover both high-level and low-level architectural concepts with a primer on the basics of web applications and the transfer of information over the internet.

\subsection{Architecture styles}

An architecture style represents a combination of design decisions and principles to organize an application in a certain way. It is the highest level of organization, not used to solve a specific problem but to enhance certain good practices in a particular type of system \cite{ghidersa_software_2022}.

\subsubsection{Monolithic Architecture}

Monolithic architecture is a software development model where all components are interdependent and interconnected; every component is required to run the application. It treats the entire code as a single program, building a new project, applying frameworks, scripts, and templates, which makes testing easier.

The main drawback to this style of architecture is that it becomes difficult to manage or update the code base as it grows. As each element is interdependent, scaling the application can be complex. Introducing even a single point of failure can break the application.

Monolithic architecture is best used for simple applications that have any future scalability planned out in advance. Knowing as much about the future scope of a project's features in advance as possible helps determine what will be feasible or not under a given structure \cite{william_web_2025}.

\subsubsection{Microservices Architecture}

Microservices architecture serves as the logical counterpoint to a monolithic environment. Using it, code is organized among independent services that don't rely on one another to perform their specific function and often communicate via RESTful APIs. Each microservice contains its own database and operates under its own business logic. Because it’s loosely coupled (meaning components do not have to change significantly if another one within the overall software does), microservice architecture makes each individual service much easier to add features to while maintaining continuous delivery. Developers can thus quickly adapt the software to changing circumstances. 

However, maintaining a piece of software with multiple services comes with drawbacks. When the number of services grows, the complexity of managing them grows, too. Microservices architecture is a good choice for highly complex applications with the resources to maintain and scale multiple independent services \cite{william_web_2025}.

\subsection{Architecture patterns}

In contrast to styles, architectural patterns are decisions and principles intended to avoid specific issues with the low-level implementation of the software. Essentially, the style comes first and determines the project's broadest level of structure, and the patterns determine the specifics of how each of the components interact with each other \cite{ghidersa_software_2022}. Below are the most prominent architectural styles and how each covers different possible use cases with their strengths and weaknesses.

\subsubsection{Preface: The Client-Server model}

The Client-Server model delineates two types of entity: clients that seek information and send structured requests for it, and servers that carry out these requests by fetching and processing the information, thus serving the clients. Often communicating through APIs, it is the server's job to stay open and available for requests, and the client's job to behave in a manner the server expects of them to best facilitate the transfer of information and handle the presentation of data, completing the path of information between the server and the end user.

Though technically a form or pattern of architecture, the idea of the client-server model is broad enough that it encompasses all web development to some extent -- for a user to get any information from the internet, there must be some sort of server they get it from. Thus, the model acts as a kind of superset comprised by the different patterns detailed below -- each is structured around clients and servers in some way, but to different extents and sometimes with very different strategies.

The model also appears throughout the field of software development at large, including other subjects pertaining to web development. In \autoref{databases}, we will discuss how certain database systems create their own client-server interface within a web server to facilitate the movement of information through the system before it is delivered to the web client. the following architectural patterns in this section are not so broad and will pertain for more directly to web architecture.

\subsubsection{Layered}

The layered pattern is a simple way to organize services in a web application. In it, information simply flows through different levels of the software to and from the user and the system in sequence, interacting with the code at each level before being passed on to the next one. The most common permutation of this pattern comprises four layers:
\begin{itemize}
	\item \emph{presentation (or UI)}, consisting of everything the user sees and interacts with. Pages, menus, text, images, etc.
	\item \emph{business}, in which the business logic of the application is defined and operated
	\item \emph{persistence}, containing mapping functions and other operations 
	\item \emph{database}, where information is persistently stored
\end{itemize}

\begin{figure}[!ht]
	\begin{center}
		\woopic{layered}{.25}
		%\vspace{-.2 in}
		\caption{An example of layered architecture}\label{layered}
	\end{center}
\end{figure}

While the simplicity of this pattern makes it easy to understand and useful for lightweight, small-scale applications, the lack of flexibility brought on by sending information through many consecutive layers can be inefficient for large applications. Since the layers are interdependent, it also introduces many potential points of failure should one or more layers develop a fault.

Another important permutation of layered architecture is the \textbf{Model-View-Controller (MVC)} design. Put simply, the MVC pattern separates the components into three roles: backend logic, frontend presentation, and an intermediary to connect the two, respectively.
\begin{itemize}
	\item \textbf{Model} -- contains business logic, core functionality, and persistent data storage. It does not contain any logic that shapes how the data is presented to the user; that is the controller’s job
	\item \textbf{View} -- Displays information to and interacts with the user. Takes data from the Model and Controller, processes it, and dynamically renders it for a user.
	\item \textbf{Controller} -- Processes the inputs the view receives from the user to send to the model and returns data to the view and the user from the model.
\end{itemize}

Broadly, it can be seen as organizing components into the frontend/backend/middleware trichotomy described earlier. This variation is less straightforward than some of the other architectural patterns, which can make it unsuitable for some projects that aren't prepared to handle its particular style of decoupling. It also illustrates how multiple functions can be performed by a single layer in such a way that the flow of data is not interrupted. This manifests in how the controller handles the input/output of both the view and model, and how the model acts both processing and data storage. MVC doesn't necessarily differ from layered architecture, but rather separates the layers into three distinct groups. The concept is illustrated in \autoref{mvc}

\begin{figure}[!ht]
\begin{center}
\woopic{MVC}{.9}
%\vspace{-.2 in}
\caption{A simple diagram of the Model-View-Controller architectural pattern}\label{mvc}
\end{center}
\end{figure}

\subsubsection{Command and Query Responsibility Segregation (CQRS)}

The signature aspect of the CQRS pattern is that it pares software down to just two main components: a user-interface and a data store. The user, through the UI, performs operations on the data store directly without having to go through an intermediary like the controller in the MVC pattern. In CQRS, data operations are separated into reading- and writing-based actions: queries and commands. Queries are used to gather data from the database storage; the user specifies what data they want and the application fetches it for them. Commands are responsible for changing data, and can carry out a wide range of actions on it. To actively modify, rather than just read, data from the store might be a more intensive task than the UI can perform on its own, which is where Command Handlers come in \cite{ghidersa_software_2022} -- if needed, these serve as a single-purpose intermediary between the two main components and have the ability and authority to make changes to the data store.

\begin{figure}[!ht]
	\begin{center}
		\woopic{CQRS}{1.0}
		%\vspace{-.2 in}
		\caption{A simple diagram of the Command \& Query Relational Segregation architectural pattern}\label{cqrs}
	\end{center}
\end{figure}

Direct access to an application's data store can be a lot of freedom to give a user. Depending on the specific data storage component the application uses, certain restrictions can be placed on what type of actions users can perform. These restrictions can essentially function as a stripped-down form of business logic for the application without a separate component to handle it. Queries and commands can either succeed or fail, depending on their compatibility with the data store's restrictions, and send the result back to the UI, whether directly or through a Command Handler.

The advantage of CQRS that it separates the intermediate layer of the application into two wholly distinct components that can operate independently -- this way more attention or resources can easily be given to one or the other if a particular project is more reading- or writing-based. It should not, however, be seen as a simple default for database-centric applications; if all you need in an application are the basic commands of creating, reading, uploading, or deleting data (collectively known as CRUD commands), CQRS is likely far more technically complex than is necessary. Usually working with GraphQL, CQRS is a pattern that is geared toward working with it and other query languages to fetch and manipulate data.

\subsubsection{Microfrontends}

Where the user-facing functionality of an application is contained in a single layer in  many other  patterns, microfrontends refer to applications that split the user interface (UI) into multiple services, usually worked on by different teams, each with different tech stacks. The pattern essentially extends the microservices style to the frontend, greatly enhancing the capability of the UI by breaking up the monolith. Instead of one single component being wholly responsible for rendering and managing the entire user interface, splitting this responsibility among microfrontends allows functions to be addressed by and linked between different components \cite{jackson_micro_2019}. For example, a site that displays a list of options for the user to choose from like a restaurant might have one frontend handle the options, formatting them and how they're organized, while another handles background aspects of the page like the broad structure and how access to other parts of the site are displayed. An example of this is shown in \autoref{microex}. 

\begin{figure}[!ht]
	\begin{center}
		\woopic{microex}{0.525}
		%\vspace{-.2 in}
		\caption{An example of a single webpage using microfrontends to manage separate sections of the site}\label{microex}
	\end{center}
\end{figure}

The key advantage of microfrontends is the ability to break the development of each frontend away from the others, allowing different teams to work independently on their designated frontend while coming into conflict with the others as little as possible. While the restaurant example is likely simple enough not to require multiple dedicated teams to maintain, a more complex application like a retail website or streaming service could conceivably have many teams with different priorities -- some might be aiming to constantly innovate and change their area of the service to keep things fresh and up to date, while others may be content simply to make sure their frontends don't break or otherwise fall behind the progress of the rest of the service.

The main disadvantage of microfrontends is in line with the potential drawbacks of microservices in general: keeping multiple disparate components interconnected while maintaining a level of independence among them can easily be an extremely intensive task. Each component must maintain a relatively consistent system of inputs and outputs that let the others easily interact with. Testing is also a potential difficulty, as running tests on individual components may not capture issues that arise when connected to the larger whole. Thus, the level of independence each frontend can maintain is lessened the more they have to interact with each other to run tests on potential changes.

The emergence over time of the microfrontend pattern exemplifies the growing importance and complexity of user interfaces in the modern web. While naturally better-suited to large development teams and resource-intensive frontends, microfrontends serve as a proof-of-concept that it's just as possible to prioritize and expand the functionality of the user interface as with any other segment of an application.

\subsubsection{JAMstack}

Where microfrontends use multiple services to make the user interface more layered and complex, JAMstack (JAM stands for "Java, API, and Markup") moves in the opposite direction. The principal attribute of JAMstack is the lack of a traditional backend server -- pages are instead sent to the user through a Content Delivery Network (CDN), which can send files to the user but does not hold data or process requests. To accommodate this, rather than use dynamic webpages as all previous patterns generally do, JAMstack instead relies solely on static, pre-rendered webpages written in HTML and JavaScript. Without a web server, the client-side JavaScript code is responsible for handling communication with the CDN and any other microservices, such as data storage or identity management.

Unlike patterns focused on optimizing for the dynamism of their frontend and/or the scalability of their backend features, JAMstack optimizes for space and cost management by eschewing dynamic frontends and many common backend features entirely. This pattern is most useful for simple web tools and other applications without the need to increase in scale by adding new features over time.

\begin{figure}[!ht]
	\begin{center}
		\woopic{jamstack}{1.1}
		%\vspace{-.2 in}
		\caption[Contrasting the JAMstack with traditional layered architecture]{Contrasting the JAMstack (right) with a traditional layered architecture (left)}\label{jamstack}
	\end{center}
\end{figure}

\subsubsection{Flux}

Flux is a pattern created by Facebook, designed to work with the service's React frontend components to build client-side web applications. It was created as an alternative to the Model-View-Controller pattern described earlier to avoid unnecessarily tangled connections between the services in an application \cite{tay_depth_2020}. Flux has a unidirectional data flow, meaning data is always sent through the services in the same sequential order. In the most basic example, the app has three services that process an action taken by a user:
\begin{itemize}
	\item \textbf{Dispatcher} -- Receives the action taken by the user and determines both where to send it, and the order information is sent in
	\item \textbf{Store} -- Contains the business logic, which it operates on data received from the dispatcher. Also manages the state of the application
	\item \textbf{View} -- Receives specifications from the store and renders the information to be presented to the user. The view re-renders when the store announces changes and is responsible for creating new actions based on user input and interaction
\end{itemize}
Each of these services only sends data in one direction and receives it from the other without a need to double back to confirm or update data. This results in a simpler application that is easier to follow and debug while each service can be optimized for its narrow focus.

Flux was designed with a specific use case in mind: managing React components simply and efficiently. Its lack of two-way inter-service communication would not be a good fit for a project managing large amounts of information, but for a user interface-focused application that doesn't send data more complex than basic text and/or images, it makes perfect sense.

\begin{figure}[!ht]
	\begin{center}
		\woopic{flux}{.3}
		%\vspace{-.2 in}
		\caption{A simple diagram of the Flux architectural pattern}\label{flux}
	\end{center}
\end{figure}

% FIT THESE IN WHEREVER THEY END UP

\section{Containerization and deployment}

The deployment of a web app is the process of transferring your software from the local machines used to develop and test it to a hosting provider that can make it accessible to end users. This usually involves paying for backend servers (though you can run them yourself) and buying a domain name to allow users to access your app as a website.

When applications move from the development to the deployment stage, the components of an application must have concrete ways of interacting with each other and the hardware (usually servers) required to become accessible to end users. A popular solution in recent times is containerization \cite{adedeji_full-stack_2023}.

Containerization is a way of packaging an application together with everything it needs to run, like its code, system libraries, and configurations, into a single unit that can run on a physical or a virtual machine (VM). By existing at the OS level, containerization allows a developer to run multiple containers inside a single OS instance. While it reduces overhead and processing power, it also makes them much faster to start and more resource-efficient.

Tools like Docker build these containers from reproducible definitions (which they call Dockerfiles), ensuring the application runs the same way on different machines, maintaining consistency in development, testing, and deployment. Containerization tools allow teams to ship applications as portable, self-contained artifacts that can run reliably across different environments such as local servers or cloud platforms.

\section{Building a tech stack for this project}

To determine what elements of a potential tech stack would be most useful to this project, this section lays out the planned extent of its capabilities and discusses potential solutions and design choices pertaining to the architecture of the software as they relate to the concepts discussed within this chapter.

First, the user needs a way to communicate their Last.fm listening history to the application, either by uploading a file generated by a Last.fm export tool, or if the project is granted direct API access, entering their username for the application to collect the data on their profile. These initial tasks require very little of the frontend -- either or both of a text field for a username and an upload window for a file. Once the user's listening history is obtained from either an input file or the site's API, it has to first go through some initial processing to generate a list of each artist they listened to, along with how many times each was played. Once these lists are generated, the most important process begins: the scraping.

Scraping is, comparatively, a much more intensive process than the other functions of this application. Consequently, the decision on how to undertake it will be important. While the idea of offloading the scraping of a listening history onto the user making the query would certainly save resources on the server side, it would not be easy or ideal -- relying on the client's browser to do the work would necessitate running the code in JavaScript, which would deprive us of some of the key tools we've planned to use for scraping and data handling. In addition, the processing undertaken by the user's browser could last several minutes, depending on the size of their listening history; this isn't ideal, as keeping a single process running in one tab for so long would put an undue burden on the user to keep that tab loaded and running. The longer such a process would run, the greater the risk of a connection or other client-side error interrupting the process. Since running the scraping client-side has these potential drawbacks, the best solution seems to involve a more standard solution: running the processing and scraping on a web server.

Using a centralized server system for the scraping gives us significant freedom in choosing how to go about it. If data from the user is collected by a client-side UI and sent to a backend server, there are numerous ways that data could be processed and returned to the user. The processing would not be restricted to client-side tools and could thus freely use non-JavaScript coding languages, as well as libraries and modules that could not be passed directly to a browser. In addition, the lack of dependence on the client-side after the initial upload means scraping and processing can potentially run without the user needing to keep their browser tab open, allowing for the possibility of implementing an extended queue of requests if users submit histories to the app faster than it can scrape the artists within them.

The primary function of the backend, once it has finished scraping the user's artists, is to perform data analysis to generate select insights, and synthesize the data into a form that can be displayed to the user. Since the use of a persistent data store is not required for these tasks, one is not strictly necessary for the project. There are, however, potential benefits an implemented data store can provide, which is discussed in \autoref{databases}.

While we determined that the scraping and processing of data is best performed on the server-side, the next step requires consideration of whether the rendering of the final data visualization to be presented to the user should also be rendered server-side. While the specifics of the data visualization used in the project is discussed in detail in \autoref{dataviz}, suffice to say that the data the application procures for visualization will not be particularly complex and can be displayed with fairly simple charting strategies. The question faced at this stage is: would it be more useful or efficient to generate the visualizations on the server-side before sending them to the client, or sending the data to the client-side for it to be generated by the user's browser?

Client-side tools like JavaScript libraries are much more prevalent for changing the visual appearance of a site than performing complex processes or calculations, so it stands to reason that generating visualizations and integrating them into the webpage would be more likely to be a task best operated on the client-side. Indeed, JavaScript libraries exist for transforming a dataset into detailed, even interactive, visualizations, which should be more than able to display the simple data from the server. With these facts considered, the best course of action is for the backend, once it's finished scraping and processing the data, to send it back to the frontend to be generated and rendered on the client-side.

In keeping with the general format of other Last.fm-based web apps and acknowledging the very specific task the application performs, it would be ideal to create it as a single-page application. For a user, the ideal experience with the site would be to upload their listening history, wait for the process to finish (if they want), and receive the results without having to navigate to a new page. If functionality is implemented to give the user a code to come back and view their results at a later date, it stands to reason that it should also be done in such a way so that they can access that from the main page, rather than a separate page for that purpose. 

Taken all together, the flow if information through the application would go like this: 

\begin{itemize}
	\item The user loads the page and makes an upload to the frontend
	\item The frontend sends the file to the backend to extract the artists and perform the scraping
	\item Optionally, some data is cached in a data store (see \autoref{databases})
	\item The backend passes the processed data back to the frontend
	\item The data is rendered into a visualization via the UI
\end{itemize}

This is an example of layered architecture, where each component is used in sequence, and the pertinent data flows in both directions. Enough of the traditional tools like a frontend framework and a centralized server are needed that using a stripped-down architecture like JAMstack or CQRS wouldn't be feasible. For a small-scale data-driven web application, layered architecture is an ideal solution.


